% ============================================================================
% Fichier  : partie3/P03_C09_volatilite_acquisition.tex : FORN (LPI342)
% Sources : RFC 3227 ; NIST SP 800-86 ; ISO/IEC 27037. Opération TSANAGA.
% ============================================================================

\chapter{Ordre de Volatilité et Acquisition à Chaud}
\label{chap:acquisition}
\chaptermark{Volatilité et acquisition à chaud}

\epigraph{%
    « On capture d'abord ce qui s'efface le plus vite. »%
}{Principe du cours}

\niveauChapitre{Intermédiaire à expert}{Chapitre~\ref{chap:duplication}.}

\begin{boxobjectifs}
\begin{itemize}
    \item Énoncer l'ordre de volatilité.
    \item Décider quand acquérir à chaud.
    \item Capturer la mémoire vive et l'état du système.
    \item Mesurer l'empreinte de l'acquisition à chaud.
    \item Documenter un arbitrage CRO.
\end{itemize}
\end{boxobjectifs}

\bigskip

% ============================================================================
\section*{Question directrice : sur un système allumé, par quoi commencer l'acquisition ?}
% ============================================================================

Sur un système vivant, certaines traces (mémoire, connexions) disparaissent à
l'extinction, d'autres (disque) sont stables. L'ordre de volatilité fixe la
priorité. Ce chapitre l'applique.

\medskip
\textbf{Plan du chapitre.}
\begin{enumerate}
    \item L'ordre de volatilité (section~\ref{sec:c09-volatilite}).
    \item L'acquisition à chaud (section~\ref{sec:c09-chaud}).
    \item Capturer la mémoire (section~\ref{sec:c09-memoire}).
    \item L'arbitrage CRO (section~\ref{sec:c09-cro}).
\end{enumerate}

\separateur

% ============================================================================
\section{L'Ordre de Volatilité}
\label{sec:c09-volatilite}
% ============================================================================

\begin{boxprocedure}
\textbf{Acquérir selon l'ordre de volatilité (RFC 3227)}
\begin{enumerate}
    \item Mémoire vive et état (processus, connexions, utilisateurs).
    \item Fichiers temporaires, données de session.
    \item Disque (image bit à bit).
    \item Journaux distants, configurations.
    \item Supports d'archives.
\end{enumerate}
\end{boxprocedure}

% ============================================================================
\section{L'Acquisition à Chaud}
\label{sec:c09-chaud}
% ============================================================================

\begin{boxdef}
\termecle{Acquisition à chaud}~: capture de données sur un système en
fonctionnement, avant son extinction. Nécessaire pour la mémoire et l'état
vivant, elle laisse inévitablement une empreinte sur le système, qu'on
minimise et documente.
\end{boxdef}

\begin{boxalert}
\textbf{Violation typique.} Éteindre un système allumé « pour le mettre sous
scellé » sans capturer la mémoire.

\smallskip
\textbf{Règle~:} si des données volatiles sont utiles (session chiffrée
ouverte, connexions, secrets en clair), on les capture d'abord. Éteindre sans
réfléchir détruit une preuve unique.
\end{boxalert}

% ============================================================================
\section{Capturer la Mémoire}
\label{sec:c09-memoire}
% ============================================================================

On capture la mémoire avec un outil à empreinte minimale, documenté, en
écrivant sur un support externe, puis on calcule le condensat de la capture.
La mémoire révèle ce que le disque ignore~: processus, connexions, clés de
chiffrement en clair.

% ============================================================================
\section{L'Arbitrage CRO}
\label{sec:c09-cro}
% ============================================================================

\begin{boxCRO}
\textbf{L'acquisition à chaud, un arbitrage.} Toucher au vivant laisse une
empreinte~: \cro{0}{-}{+}. On accepte une légère altération (fiabilité) pour
sauver une trace unique et opposable (opposabilité), et on la documente pour
la rendre défendable. Le choix se justifie et se trace.
\end{boxCRO}

\begin{boxscenario}{Opération TSANAGA. Une session chiffrée ouverte}
\textbf{Constatation.} À la saisie, l'ordinateur de M.~NDONGO était allumé,
avec un conteneur chiffré monté (déchiffré en mémoire).

\medskip
\textbf{Analyse} (confiance élevée). Éteindre aurait reverrouillé le
conteneur, rendant l'accès très difficile. L'équipe a capturé la mémoire à
chaud (où la clé était en clair), puis acquis le disque. L'arbitrage, risqué,
a été documenté et justifié.
\end{boxscenario}

\begin{boxmetier}{Analyste forensique}
\textbf{Décider vite, tracer tout.} Sur une scène vivante, l'analyste décide
en minutes, mais documente chaque choix~: pourquoi capturer à chaud, avec
quel outil, à quelle heure. La trace de la décision vaut autant que la
capture.
\end{boxmetier}

\separateur

\begin{boxresume}
\textbf{Ce qu'il faut retenir de ce chapitre~:}
\begin{itemize}
    \item \textbf{L'ordre de volatilité}~: mémoire, puis disque, puis
          archives.
    \item \textbf{L'acquisition à chaud} sauve les traces volatiles, en
          laissant une empreinte.
    \item \textbf{On capture la mémoire} avant l'extinction, si elle est
          utile.
    \item \textbf{L'arbitrage CRO} se justifie et se documente.
    \item \textbf{Une session chiffrée ouverte} change la décision.
\end{itemize}
\end{boxresume}

\begin{boxinfo}
\textbf{Points de contrôle de l'élève.} Je sais~: énoncer l'ordre de
volatilité~; décider d'acquérir à chaud~; capturer la mémoire~; analyser le
CRO de l'acquisition~; documenter l'arbitrage.
\end{boxinfo}

\bigskip

\begin{boxexo}[title={\ding{45}\quad Exercice~9.1. Volatilité \niveaubase}]
Classez du plus volatil au plus stable~: disque, mémoire, archives.
\end{boxexo}

\begin{boxexo}[title={\ding{45}\quad Exercice~9.2. À chaud \niveaubase}]
Qu'est-ce qu'une acquisition à chaud~?
\end{boxexo}

\begin{boxexo}[title={\ding{45}\quad Exercice~9.3. Support \niveaubase}]
Sur quel support écrit-on une capture mémoire~?
\end{boxexo}

\begin{boxexo}[title={\ding{45}\quad Exercice~9.4. Éteindre \niveaumoyen}]
Pourquoi ne pas éteindre un système allumé avant de capturer la mémoire~?
\end{boxexo}

\begin{boxexo}[title={\ding{45}\quad Exercice~9.5. Conteneur chiffré \niveaumoyen}]
Pourquoi la session chiffrée ouverte change-t-elle la décision~?
\end{boxexo}

\begin{boxexo}[title={\ding{45}\quad Exercice~9.6. CRO \niveaumoyen}]
Expliquez l'arbitrage CRO de l'acquisition à chaud.
\end{boxexo}

\begin{boxexo}[title={\ding{45}\quad Exercice~9.7. Empreinte \niveaumoyen}]
Pourquoi documenter l'empreinte laissée par l'outil d'acquisition~?
\end{boxexo}

\begin{boxexo}[title={\ding{45}\quad Exercice~9.8. Ce que révèle la mémoire \niveaumoyen}]
Citez trois éléments que la mémoire révèle et que le disque ignore.
\end{boxexo}

\begin{boxexo}[title={\ding{45}\quad Exercice~9.9. Procédure \niveauexpert}]
Rédigez la procédure d'acquisition à chaud de l'ordinateur de TSANAGA.
\end{boxexo}

\begin{boxexo}[title={\ding{45}\quad Exercice~9.10. Justification \niveauexpert}]
Rédigez la note qui justifie l'acquisition à chaud devant le tribunal.
\end{boxexo}

\bigskip

\begin{boxinfo}
\textbf{Pour aller plus loin.}
\begin{itemize}
    \item RFC~3227~\cite{rfc3227}~; NIST SP~800-86~\cite{nist80086}~;
          ISO/IEC 27037~\cite{iso27037}.
\end{itemize}
\end{boxinfo}

\begin{boxinfo}
\textbf{Transition.} L'acquisition faite, l'analyse commence~: la Partie~IV
parcourt les domaines, à commencer par le disque
(Chapitre~\ref{chap:disque}).
\end{boxinfo}
